\chapter{AgentOS 有限认知界面定理}

\begin{chapterlead}
前一章已经把工作记忆 $h(t)$ 从容量较小的``短期存储区''重新定位为多尺度认知结构在当前时刻显式交汇、竞争、组合与行动闭合的稀疏认知总线，并指出成熟智能的目标不是让更多内容进入工作记忆，而是让越来越少、越来越关键的结构在正确时刻进入工作记忆。一个伴随这一结论自然产生的问题是：这样的有限界面为什么必然存在？本章将回答这个问题。我们将看到，有限认知界面不是工程实现上的权宜安排，而是由开放生命周期、有限计算资源与有限响应时间三条基本条件共同强制成立的结构性事实，因而可以表述为一条定理。它是第二编工作记忆理论线的基础定理，后续关于自我感知、认知控制与运行态治理的机制，都将以这一界面为前提展开。

论证从生命周期信息的无界增长开始。持续运行的智能体不断接收输入、形成经历并更新内部结构，其历史规模原则上没有与当前计算资源同尺度的固定上界；但任意时刻的感知、推理与行动都不可能重新处理全部历史，因此必须区分不断增长的长期状态 $S_t=(\Theta_t,E_t)$ 与当前真正参与计算的状态 $h(t)$，智能体的运行被抽象为 $(S_t,x_t)\xrightarrow{\Pi_t}h(t)\xrightarrow{F}a_t$。在给出当前认知界面的定义与有限性条件之后，本章从有限计算资源与有限响应时间两条相互独立的路径证明 $\sup_t\mathcal{C}(h(t))<\infty$：生命周期状态可以持续增长，而当前认知界面必须保持有限。在此基础上，本章相继澄清三个关键认识：这一定理限制的是 $h(t)$ 而非长期记忆，``可以被保存的信息''与``当前必须被直接处理的信息''必须在体系结构上分离；有限界面不是对最近历史的固定窗口截取，而是长期状态在当前时刻形成的有效投影；因此对超长历史的处理本质上不是把历史摘要越压越短的历史压缩，而是把历史逐步转化为情景状态与结构状态的状态压缩。由此，当前认知界面的核心问题被表述为容量受限下的信息选择问题：目标不是容纳更多信息，而是在有限状态中保留与当前任务最相关的信息。

本章最后给出定理的工程含义、它与三相记忆动力学的关系以及几何解释：运行了十分钟的智能体与运行了十年的智能体可以共用同一个有限大小的输入输出界面，无限生命周期由无限多个有限认知时刻 $h(t)\to a_t\to h(t+1)$ 连接而成，$h(t)$ 是整个生命周期记忆结构 $\mathcal{M}$ 在当前时刻形成的有限截面；并由此得到三条直接推论：长期记忆规模与单次模型输入规模可以解耦，超长生命周期不等于超长模型上下文，真正需要扩展的是长期状态空间而不是当前上下文窗口。本章在论证上只使用已经建立的三相记忆结构与认知总线结论，不预先依赖后续的具体模块；它确立的是``界面必须有限''，而``什么应当进入 $h(t)$''的选择机制将首先在下一章 SPC：自我感知压缩器中获得一个具体实例，有限界面内部能够承载的认知复杂度边界，则留待后续进一步建立。
\end{chapterlead}

\section{从无限生命周期到有限当前状态}

一个能够长期运行的智能体，其生命周期信息量会随着时间不断增加。智能体持续接收新的世界输入，形成新的经历，更新既有结构，并产生新的内部状态。如果将从初始时刻到当前时刻的全部历史信息记为

\[
\mathcal{H}_{0:t},
\]

那么对于一个开放运行的智能体，通常有

\[
|\mathcal{H}_{0:t}|
\rightarrow \infty,
\qquad
t\rightarrow\infty.
\]

这里的无限并不要求在有限时间内实际达到数学意义上的无穷，而是表示：智能体的生命周期长度原则上没有一个与当前计算资源相同尺度的固定上界。

然而，智能体在任意一个具体时刻进行感知、推理和行动时，并不可能同时重新处理其全部生命周期历史。

这意味着必须区分两个不同层面的量：

\[
\text{生命周期状态}
\]

与

\[
\text{当前认知状态}.
\]

在三相记忆结构中，可以将智能体在时刻 $t$ 已形成的主要内部状态表示为

\[
S_t
=
\left(
\Theta_t,
E_t
\right),
\]

其中 $\Theta_t$ 表示经过长期形成的结构记忆，$E_t$ 表示仍然以经历形式保存的情景记忆。

外部世界在当前时刻产生输入

\[
x_t,
\]

而智能体真正用于当前计算的状态记为

\[
h(t).
\]

因此，一个长期运行的智能体并不是直接建立如下关系：

\[
(\mathcal{H}_{0:t},x_t)
\rightarrow
a_t,
\]

而是必须首先形成一个当前状态

\[
h(t),
\]

再由当前状态产生动作：

\[
h(t)
\rightarrow
a_t.
\]

由此，智能体的运行可以抽象为

\[
(S_t,x_t)
\xrightarrow{\Pi_t}
h(t)
\xrightarrow{F}
a_t,
\]

其中 $\Pi_t$ 表示从长期状态和当前输入形成当前认知状态的映射，$F$ 表示在当前状态上完成推理和行动生成的计算过程。

这里的 $h(t)$ 就构成了 AgentOS 的\textbf{当前认知界面}。


\section{当前认知界面的定义}

\subsection{定义}

\textbf{定义：当前认知界面}

对于一个在时刻 $t$ 运行的智能体，如果存在状态

\[
h(t),
\]

使得当前时刻所有直接参与感知、推理和行动生成的有效信息，都必须通过 $h(t)$ 进入当前计算过程，则称 $h(t)$ 为该智能体在时刻 $t$ 的当前认知界面。

因此有

\[
a_t
=
F(h(t)).
\]

当前认知界面不是智能体全部记忆的集合，也不是全部生命周期历史的副本，而是智能体在某一具体时刻能够直接用于当前计算的状态切面。

因此应当区分

\[
S_t
\]

和

\[
h(t).
\]

前者描述智能体已经形成的长期状态，后者描述当前正在被直接激活和使用的状态。

一般而言，

\[
h(t)
\subsetneq S_t \cup \{x_t\},
\]

这里的包含关系并不表示简单的数据集合包含，而表示 $h(t)$ 所包含的信息来源于长期状态和当前世界输入，但只保留当前计算实际需要的部分。


\subsection{有限认知界面}

如果存在一个与智能体生命周期长度无关的有限常数 $C_h$，使得任意时刻都有

\[
\mathcal{C}(h(t))
\leq
C_h,
\]

则称该智能体具有\textbf{有限认知界面}。

其中

\[
\mathcal{C}(h(t))
\]

表示当前认知界面的规模。

这里的规模不能仅理解为文本 token 数。根据具体实现，它可以表示：

\[
\mathcal{C}(h(t))
=
\text{当前直接参与计算的信息规模},
\]

例如可以包括：

\begin{itemize}
    \item 当前输入的信息量；
    \item 当前激活的记忆量；
    \item 当前状态变量的数量；
    \item 当前推理所直接依赖的信息结构。
\end{itemize}

因此，有限认知界面首先是一个计算约束，而不是一种特定的大模型上下文格式。


\section{有限认知界面定理}

\subsection{定理}

\textbf{定理：AgentOS 有限认知界面定理}

设一个智能体满足以下条件：

\begin{enumerate}
    \item 智能体具有开放的生命周期，并能够持续获得新的非完全冗余信息；
    \item 智能体在任意单个运行时刻所拥有的直接计算资源有限；
    \item 智能体需要在有限时间内持续地产生新的状态或动作；
    \item 生命周期历史可以通过内部状态保存，而不要求全部历史信息同时直接参与当前计算。
\end{enumerate}

则对于任意能够持续运行的此类智能体，必然存在一个当前认知状态

\[
h(t),
\]

以及从长期状态和当前输入到该状态的映射

\[
\Pi_t:
(S_t,x_t)
\rightarrow
h(t),
\]

使得

\[
\mathcal{C}(h(t))
\leq
C_h,
\]

其中 $C_h$ 为与生命周期长度无关的有限上界。

也就是说，即使智能体的生命周期状态规模不断增长，

\[
\mathcal{C}(S_t)
\rightarrow
\infty,
\]

当前直接参与一次认知计算的状态仍然必须满足

\[
\sup_t
\mathcal{C}(h(t))
<
\infty.
\]

因此：

\begin{statementbox}
\centering
\adjustbox{max width=.96\linewidth}{\(\displaystyle
\text{生命周期状态可以持续增长，而当前认知界面必须保持有限。}
\)}
\end{statementbox}


\section{定理证明}

\subsection{由有限计算资源证明}

设智能体在任意一个运行周期内可用于当前计算的最大计算资源为

\[
B<\infty.
\]

设处理一个最小有效信息单位至少需要非零计算代价

\[
c>0.
\]

那么一次运行周期内可以直接参与计算的信息单位数量至多为

\[
N_{\max}
\leq
\frac{B}{c}.
\]

因此，当前实际参与计算的状态规模必然满足

\[
\mathcal{C}(h(t))
\leq
\frac{B}{c}.
\]

令

\[
C_h
=
\frac{B}{c},
\]

则有

\[
\mathcal{C}(h(t))
\leq
C_h.
\]

另一方面，如果智能体持续获得新的有效信息，则其生命周期状态规模可以随着时间持续增加：

\[
\mathcal{C}(S_t)
\rightarrow
\infty.
\]

假设不存在有限认知界面，而要求每一次当前计算都直接使用完整生命周期状态，即

\[
h(t)=S_t,
\]

则必然有

\[
\mathcal{C}(h(t))
=
\mathcal{C}(S_t).
\]

由于

\[
\mathcal{C}(S_t)
\rightarrow
\infty,
\]

因此总存在某一时刻 $t^\star$，使得

\[
\mathcal{C}(S_{t^\star})
>
C_h.
\]

此时完整生命周期状态所需的计算资源已经超过单周期可用资源，智能体无法按照假设继续完成当前计算。

这与智能体能够持续运行的条件矛盾。

因此必然存在一个从长期状态到有限当前状态的映射

\[
\Pi_t:
(S_t,x_t)
\rightarrow
h(t),
\]

并满足

\[
\mathcal{C}(h(t))
\leq
C_h.
\]

定理得证。


\subsection{由有限响应时间证明}

上述结论也可以从时间约束得到。

设智能体当前需要直接处理的信息量为

\[
L_t,
\]

处理单位信息至少需要时间

\[
\tau_0>0.
\]

则完整处理这些信息所需时间至少满足

\[
T_t
\geq
\tau_0 L_t.
\]

如果要求每一次决策都完整读取全部生命周期历史，则随着生命周期增长，

\[
L_t
\rightarrow
\infty,
\]

从而有

\[
T_t
\rightarrow
\infty.
\]

然而，一个与外部世界持续交互的智能体必须在有限时间内产生下一状态或动作，因此存在有限响应上界

\[
T_{\max}<\infty.
\]

要满足

\[
T_t
\leq
T_{\max},
\]

则必须有

\[
L_t
\leq
\frac{T_{\max}}{\tau_0}.
\]

因此当前直接参与计算的信息量存在有限上界。

这再次说明：

\[
\boxed{
\mathcal{C}(h(t))
\leq
C_h.
}
\]


\section{长期记忆与当前认知界面的分离}

有限认知界面定理并不意味着智能体只能拥有有限的长期记忆。

它限制的是

\[
h(t),
\]

而不是

\[
\Theta_t
\]

或者

\[
E_t.
\]

因此完全可以出现：

\[
|\Theta_t|
\uparrow,
\]

\[
|E_t|
\uparrow,
\]

同时保持

\[
|h(t)|
\leq
C_h.
\]

也就是说，AgentOS 必须在体系结构上区分：

\[
\text{可以被保存的信息}
\]

与

\[
\text{当前必须被直接处理的信息}.
\]

前者构成长期记忆空间，后者构成当前认知界面。

因此：

\begin{statementbox}
\centering
\adjustbox{max width=.96\linewidth}{\(\displaystyle
\text{长期记忆容量}
\neq
\text{当前认知容量}.
\)}
\end{statementbox}

随着生命周期不断延长，二者之间的比例甚至可以满足

\[
\frac{\mathcal{C}(h(t))}
{\mathcal{C}(S_t)}
\rightarrow
0.
\]

这并不意味着智能体逐渐失去过去，而意味着越来越多的过去以非当前激活状态存在。


\section{有限认知界面不是固定历史窗口}

需要特别区分有限认知界面与传统固定长度上下文窗口。

一种最简单的有限历史窗口可以写成

\[
h(t)
=
\{x_{t-n},\ldots,x_t\}.
\]

这种方法仅仅保留距离当前最近的一段历史。

但对于长期智能体而言，当前最重要的信息并不一定是时间上最近的信息。

某一当前状态可能同时依赖：

\begin{itemize}
    \item 当前世界输入；
    \item 数秒前发生的局部事件；
    \item 数日前形成的一段经历；
    \item 很久以前形成的重要关系；
    \item 已经长期稳定存在的结构记忆。
\end{itemize}

因此 AgentOS 中的当前认知界面不应写成简单的时间窗口，而应写成

\[
h(t)
=
\Pi_t
\left(
\Theta_t,
E_t,
x_t
\right).
\]

这意味着 $h(t)$ 是长期状态在当前时刻形成的有效投影。

因此：

\begin{statementbox}
\centering
\adjustbox{max width=.96\linewidth}{\(\displaystyle
h(t)
\text{ 是当前状态投影，而不是历史尾部截取。}
\)}
\end{statementbox}


\section{从历史压缩到状态压缩}

传统的长上下文压缩通常可以表示为

\[
\mathcal{H}_{0:t}
\rightarrow
\text{Summary}_t.
\]

其基本目标仍然是将较长的文本历史压缩为较短的文本历史。

AgentOS 的目标则不同。

随着经历积累，过去的信息可以逐渐由具体事件转化为更加稳定的状态。

例如：

\[
e_1,e_2,\ldots,e_n
\]

可以逐渐形成某种结构状态

\[
\theta.
\]

因此系统经历的是

\[
\text{历史信息}
\rightarrow
\text{情景状态}
\rightarrow
\text{结构状态}.
\]

在当前时刻，系统不必重新读取所有形成这些状态的原始历史，只需要使用已经形成的有效状态。

于是：

\[
\mathcal{H}_{0:t}
\not\rightarrow
h(t)
\]

而是：

\[
\mathcal{H}_{0:t}
\rightarrow
(\Theta_t,E_t)
\rightarrow
h(t).
\]

这意味着 AgentOS 对超长历史的处理，本质上不是不断生成越来越长的历史摘要，而是持续把历史转化为可以被未来重新使用的状态。


\section{当前认知界面的信息选择问题}

有限认知界面意味着一个新的问题：

如果

\[
\mathcal{C}(S_t)
\gg
C_h,
\]

那么当前状态 $h(t)$ 应当保留什么？

显然，目标不能是让 $h(t)$ 包含尽可能多的信息。

设当前任务真正相关的目标变量为

\[
Y_t.
\]

理想的当前认知界面应当在有限容量下尽可能保留与 $Y_t$ 有关的信息。

因此可以将当前认知界面的理想形式写成

\[
\Pi_t^\star
=
\arg\max_{\Pi_t}
I(h(t);Y_t),
\]

满足

\[
h(t)
=
\Pi_t(S_t,x_t)
\]

以及

\[
\mathcal{C}(h(t))
\leq
C_h.
\]

其中

\[
I(h(t);Y_t)
\]

表示当前认知状态与当前任务所需信息之间的互信息。

这说明有限认知界面的根本问题不是

\[
\text{如何容纳更多信息},
\]

而是

\begin{statementbox}
\centering
\adjustbox{max width=.96\linewidth}{\(\displaystyle
\text{如何在有限当前状态中保留最有用的信息}.
\)}
\end{statementbox}

因此，当前认知界面本质上是一个受限的信息选择问题。


\section{固定大小输入输出界面的工程含义}

有限认知界面定理进一步给出了 AgentOS 的一个重要工程原则。

假设底层推理模型每一次运行允许的最大输入容量为

\[
C_{\mathrm{model}}.
\]

AgentOS 并不需要使

\[
C_{\mathrm{model}}
\]

随生命周期不断增加。

相反，可以始终维持

\[
\mathcal{C}(h(t))
\leq
C_{\mathrm{model}}.
\]

因此，从大模型的角度看，一个运行了十分钟的 Agent 与一个运行了十年的 Agent，都可以通过同一种有限大小的输入接口完成当前推理。

二者的区别并不主要发生在当前界面的大小，而发生在界面背后的长期状态：

\[
(\Theta_t,E_t).
\]

由此得到：

\begin{statementbox}
\centering
\adjustbox{max width=.96\linewidth}{\(\displaystyle
\text{Memory grows, while the cognitive interface remains bounded.}
\)}
\end{statementbox}

即：

\begin{statementbox}
\centering
\adjustbox{max width=.96\linewidth}{\(\displaystyle
\text{记忆可以增长，而当前认知界面不必随之增长。}
\)}
\end{statementbox}

这正是 AgentOS 能够支持超长生命周期运行的基础条件。


\section{输入界面与输出界面的闭合}

从运行结构上看，AgentOS 每一个时刻都可以被视为一个有限输入输出系统。

当前输入界面为

\[
h(t),
\]

推理过程为

\[
F,
\]

输出为

\[
a_t.
\]

因此：

\[
h(t)
\xrightarrow{F}
a_t.
\]

动作作用于环境后产生新的输入：

\[
a_t
\rightarrow
W_{t+1}
\rightarrow
x_{t+1}.
\]

与此同时，本轮经历又可以更新长期状态：

\[
(\Theta_t,E_t)
\rightarrow
(\Theta_{t+1},E_{t+1}).
\]

于是整个生命周期可以表示为：

\[
(\Theta_t,E_t,x_t)
\xrightarrow{\Pi_t}
h(t)
\xrightarrow{F}
a_t
\xrightarrow{W}
x_{t+1},
\]

并进一步得到

\[
(\Theta_{t+1},E_{t+1},x_{t+1})
\xrightarrow{\Pi_{t+1}}
h(t+1).
\]

因此，虽然智能体的生命周期可以无限延伸，但每一个运行时刻都只需要完成一个有限状态转换：

\[
\boxed{
h(t)
\rightarrow
a_t
\rightarrow
h(t+1).
}
\]

无限生命周期由无限多个有限认知时刻连接而成，而不是由某一个无限大的上下文承担。


\section{与三相记忆动力学的关系}

有限认知界面定理与三相记忆结构具有直接关系。

工作记忆

\[
h(t)
\]

对应当前正在直接参与认知活动的状态。

情景记忆

\[
E_t
\]

保存已经离开当前状态、但仍然保持经历结构的信息。

结构记忆

\[
\Theta_t
\]

保存由长期经历逐渐形成的稳定结构。

因此，生命周期信息并不需要永久驻留在 $h(t)$ 中。

随着时间演化，可以发生

\[
h(t)
\rightarrow
E
\rightarrow
\Theta.
\]

而当前需要的过去信息，则可以重新影响新的

\[
h(t').
\]

因此有限认知界面不是对长期记忆的否定，而是三相记忆分层存在的必要条件。

如果所有历史信息都必须永久同时存在于 $h(t)$ 中，那么

\[
E
\]

和

\[
\Theta
\]

作为不同记忆相态将失去独立意义。

正因为

\[
h(t)
\]

保持有限，过去才必须逐渐离开当前直接激活状态，并形成不同时间尺度的记忆结构。


\section{当前认知界面的几何解释}

在 AgentOS 的生命周期记忆模型中，可以将整个生命周期记忆结构记为

\[
\mathcal{M}.
\]

随着时间推进，$\mathcal{M}$ 不断积累新的结构。

当前工作记忆

\[
h(t)
\]

并不是整个 $\mathcal{M}$，而可以理解为生命周期记忆结构在当前时刻形成的一个有限截面：

\[
h(t)
=
\Sigma_t(\mathcal{M},x_t),
\]

其中 $\Sigma_t$ 表示当前时刻的状态截面。

因此，随着生命周期延长，

\[
\mathcal{M}
\]

可以不断扩大，而

\[
h(t)
\]

仍然保持有限。

从这一视角看，智能体的``现在''不是全部过去历史的简单累加，而是整个生命周期状态与当前世界发生作用时形成的有限切面。

因此：

\begin{statementbox}
\centering
\adjustbox{max width=.96\linewidth}{\(\displaystyle
\text{过去可以无限延伸，而现在必须保持有限。}
\)}
\end{statementbox}


\section{有限认知界面定理的几个直接推论}

\subsection{推论一：长期记忆规模与单次模型输入规模可以解耦}

如果

\[
\mathcal{C}(h(t))
\leq
C_h,
\]

那么即使长期记忆满足

\[
\mathcal{C}(S_t)
\gg
C_h,
\]

底层推理模型仍然只需要处理 $h(t)$。

因此：

\[
\frac{\partial \mathcal{C}(h(t))}
{\partial \mathcal{C}(S_t)}
\]

不要求接近 $1$。

理想情况下，随着长期记忆持续增加，单次认知界面仍然可以保持同一数量级。


\subsection{推论二：超长生命周期不等于超长模型上下文}

如果一个 Agent 已经运行极长时间，其生命周期历史可能远远超过任意模型上下文窗口。

但只要存在

\[
(\Theta_t,E_t)
\rightarrow
h(t)
\]

的状态形成过程，则模型并不需要直接看到全部历史。

因此：

\begin{statementbox}
\centering
\adjustbox{max width=.96\linewidth}{\(\displaystyle
\text{Long Lifetime}
\neq
\text{Long Context Window}.
\)}
\end{statementbox}


\subsection{推论三：真正需要扩展的是长期状态空间}

当智能体能力增长时，最直接需要扩展的并不必然是 $h(t)$ 的容量。

更重要的是：

\[
\Theta_t,
\qquad
E_t,
\]

以及长期状态到当前状态之间映射

\[
\Pi_t
\]

的质量。

因此，AgentOS 的长期扩展方向可以表示为：

\begin{statementbox}
\centering
\adjustbox{max width=.96\linewidth}{\(\displaystyle
\text{扩大可保存、可形成、可重新利用的长期状态空间，}
\)}
\end{statementbox}

而不是简单地使当前上下文窗口无限增长。


\section{与下一定理的关系}

有限认知界面定理首先回答的是：

\begin{statementbox}
\centering
\adjustbox{max width=.96\linewidth}{\(\displaystyle
\text{为什么长期智能体必须存在一个有限的当前认知界面？}
\)}
\end{statementbox}

它从生命周期增长、有限计算资源和有限响应时间三个基本条件出发，证明了

\[
h(t)
\]

不能随着全部生命周期历史无界增长。

但这一结论还没有回答另一个更加具体的问题：

即使

\[
h(t)
\]

在规模上保持有限，随着同时进入当前认知界面的状态、关系和依赖不断增加，其内部计算复杂度是否仍然可能失控？

这需要进一步讨论当前工作记忆本身的复杂度边界。

因此，在有限认知界面定理之后，可以进一步建立工作记忆复杂度的有限性条件。

前者约束的是：

\begin{statementbox}
\centering
\adjustbox{max width=.96\linewidth}{\(\displaystyle
\text{当前认知面的总体规模}
\)}
\end{statementbox}

后者进一步研究：

\begin{statementbox}
\centering
\adjustbox{max width=.96\linewidth}{\(\displaystyle
\text{有限认知面内部能够同时承载多大的认知复杂度}.
\)}
\end{statementbox}

两者共同构成 AgentOS 工作记忆理论的基础约束。
